\chapter{角運動量、スピン、統計}

\section{通常の回転から内在的スピンへ}

角運動量が形式的な主題であるが、真の標的はスピンである。スピンは素粒子の最も単純かつ最も直観に反する性質の一つであり、その構造は角運動量の数学から導かれる。

角運動量はベクトルである。その向きは回転軸に沿っている。その軸に沿って二つの可能な向きがあるため、どちらを選ぶかを決める規約が必要である。右手の指を回転の向きに曲げると、親指が角運動量ベクトルの向きを指す。これは定義であり、定理ではない。

その大きさは、どれだけの質量が存在するか、その質量がどのように分布しているか、系の大きさ、およびその角速度に依存する。質量分布と大きさは慣性モーメントにまとめられる。したがって、大きな物体と小さな物体は、非常に異なる速さで回転し、非常に異なるエネルギーを持ちながら、同じ角運動量を担うことができる。

通常の複合物体は、概念的に異なる二種類の角運動量を担うことができる。外部軸の周りを円運動するカップは、たとえカップ自身を通るどの軸の周りでも回転していなくても、その重心がその軸の周りを運動しているため角運動量を持つ。これが軌道角運動量である。スピンとは内部軸の周りの回転を意味する。

明確な定義は重心系を用いる。重心が静止している系に移ると、通常の全運動量は消える。その系に残る角運動量があれば、それがスピン角運動量である。物体は内部的に静止している必要はなく、その重心だけが静止している。

バスケットボールにとって、この区別は何の困難も引き起こさない。静止したバスケットボールと、回転させた同じバスケットボールは、同じ物体が二つの異なる状態にあるものとして明らかに認識される。量子力学は、角運動量が離散的であるため、この描像を複雑にする。角運動量ゼロの状態と最初の回転状態とを結ぶ、許された角運動量の連続的な系列は存在しない。したがって、最初の回転状態が運動している同じ物体なのか、それともまったく別の物体なのかを問うことができる。

その表現は一部は慣習的であるが、物理的な問いは正確である。最初の非ゼロ回転状態に到達するにはどれだけのエネルギーが必要か。ギャップが小さければ、励起された物体は容易に認識できるままである。ギャップが莫大であれば、励起は物体を根本的に変えうる。原子に過剰な回転励起を与えようとすると、同じ原子が静かに回転しているままになるのではなく、電子をはぎ取り原子を破壊しうる。

次に電子を考え、その既知の内在的スピンはしばらく脇に置く。電子を内部軸の周りに回転させることができるだろうか。数学的な点には回転させるべき目に見える内部構造はないが、実験は、回っているのを観察できる小さな剛体表面を与えることによって電子の大きさを確定するわけではない。電子が極めて小さな有限の大きさを持つならば、その最初の追加的な回転励起は、単に到達するのに大きすぎるエネルギーを必要とするかもしれない。エネルギーは大きさが減少するにつれて増大するため、必要なエネルギーは天文学的なもの、おそらくはプランクエネルギー程度にさえなりうる。この見積もりは意図的に暫定的である。\footnote{この見積もりの直後に続く短い一節は、さらなる実質的な再構成を支えるのに十分なほど理解可能ではない。}

素粒子にとって、スピンの大きさはしたがって固定された特性として扱われる。電子は、次第に大きくなる内在的スピンの大きさの範囲を伴って存在するわけではない。その他の点では電子に似た励起が異なる全スピンを担うならば、それは異なる粒子として分類されるであろう。

固定された大きさは向きを固定しない。\emph{スピントロン}と呼ばれる粒子を想像しよう。これは部分的にはその角運動量の大きさによって同定される。そのスピンがある軸に沿って向いているならば、回転不変性は、同じ種類の粒子がそのスピンを別の軸に沿って向けることもできることを要求する。空間的な方向で本質的に優先されるものはない。スピンの量は固定されたままであるが、その向きは変化しうる。したがって、この問題には一見矛盾する二つの要素がある。量子化された大きさと、自由に回転できる向きである。角運動量代数が、それらがどのように共存するかを示すであろう。

\begin{classroomqa}
\lecturetimestamp{00:12:25}
\audiencequestion{なぜ小さな粒子を回転させるのにエネルギーが必要なのか。}
\lecturerresponse{質量\(M\)、特徴的半径\(R\)の物体について、慣性モーメントは
\begin{equation}
I\sim cMR^2,
\end{equation}
の程度である。ここで数値因子\(c\)は形状と質量の分布の仕方に依存する。外側付近の物質が速さ\(v\)で運動するなら、その通常の運動量は\(Mv\)の程度であり、角運動量は運動量×距離である：
\begin{equation}
L\sim MvR.
\end{equation}
等価に、\(L=I\omega\)で\(v\sim\omega R\)である。回転運動エネルギーは
\begin{equation}
E_{\mathrm{rot}}=\frac{L^2}{2I}.
\end{equation}
\(M\)と\(L\)を固定すると、
\begin{equation}
E_{\mathrm{rot}}\sim\frac{L^2}{2cMR^2}\propto\frac{1}{R^2}.
\end{equation}
物体が小さいほど、与えられた角運動量に要するエネルギーは大きい。この逆二乗依存性は古典的である。追加の量子力学的事実は、角運動量が離散的な段階で変化するため、最初に許される段階を任意に小さくすることはできないということである。}
\end{classroomqa}

これから用いる数学は抽象的で、最初はかなり直観に反する。その興味は、実験的に検証可能な事実が少数の代数的規則から導かれる点にある。

\section{軌道角運動量}

原点のまわりを円軌道で運動する1個の粒子から始める。その角運動量の大きさは、運動量×原点からの垂直距離に等しい。平面内では、一方の回転の向きを正、逆の向きを負と呼んでよい。符号は、選んだ軸に対する角運動量ベクトルの向きを記録する。

これはまだスピンではない。粒子が角運動量をもつのは、通常の運動量をもつからである。次に、直径の両端に2個の粒子を置き、同じ回転の向きで運動させる。それらの通常の運動量は打ち消し合うので重心は静止しているが、それらの角運動量は同じ向きを指し、加算される。この系は全運動量がゼロで全角運動量が非ゼロである：スピンの古典模型である。

原点に対する粒子の位置を
\begin{equation}
\vec r=(x,y,z)=(x_1,x_2,x_3),
\end{equation}
とし、その運動量を
\begin{equation}
\vec p=(p_x,p_y,p_z)=(p_1,p_2,p_3).
\end{equation}
とする。角運動量はベクトル積
\begin{equation}
\vec L=\vec r\times\vec p.
\end{equation}
である。因子を逆にするとベクトルも逆になる：\(\vec p\times\vec r=-\vec r\times\vec p\)。標準的な規約は\(\vec r\times\vec p\)であり、その向きは右手の規則で定まる。\footnote{\raggedright 講義動画フレーム\nolinkurl{lecture_07_figure_02.png}、00：18：49。\par}

\begin{figure}[h]
\centering
\includegraphics[width=0.72\textwidth]{lecture_07_figure_02.png}
\par
\begin{tikzpicture}[scale=1.0]
\draw[->] (-0.6,0) -- (4.8,0) node[right] {$x$};
\draw[->] (0,-0.6) -- (0,3.8) node[above] {$y$};
\coordinate (P) at (3.2,2.3);
\draw[->,thick] (0,0) -- (P) node[midway,above left] {$\vec r$};
\draw[dashed] (P) -- (3.2,0);
\draw[dashed] (P) -- (0,2.3);
\draw[->,thick] (P) -- ++(-1.0,1.25) node[above] {$\vec p$};
\node at (4.15,2.75) {$\vec L=\vec r\times\vec p$};
\end{tikzpicture}
\caption{位置\(\vec r\)、運動量\(\vec p\)の粒子について、角運動量\(\vec L=\vec r\times\vec p\)はそれらの平面に垂直である。}
\lectureframe{00:18:49}
\end{figure}

デカルト成分では、
\begin{equation}
\begin{aligned}
L_x&=yp_z-zp_y,\qquad
L_y=zp_x-xp_z,\\
L_z&=xp_y-yp_x.
\end{aligned}
\end{equation}
最初の式を覚え、ラベルを巡回させれば十分である
\[
x\longrightarrow y\longrightarrow z\longrightarrow x.
\]
成分ラベルと右辺のラベルの両方を巡回させると、相対符号を変えずに次の式が得られる。この記憶法が有用なのは、まさに符号を誤って逆にしやすいからである。

点粒子について、\(\vec r\times\vec p\)は選んだ原点のまわりの軌道角運動量である。複合系については、すべての構成要素の角運動量をベクトルとして加える：
\begin{equation}
\vec L_{\mathrm{tot}}=\sum_a \vec r_a\times\vec p_a.
\end{equation}
この和を重心系で評価すると、
\[
\sum_a\vec p_a=0.
\]
であっても非ゼロのままでありうる。その残りの角運動量が複合系のスピンである。

\section{量子演算子と正準交換関係}

量子力学では、観測可能量は演算子で表される。したがって位置、運動量、角運動量の各成分は演算子である。それらの微分演算子としての詳細な実現はここでは必要ない。交換関係に、角運動量のスペクトルを導くのに十分な情報が含まれている。

2つの演算子\(A\)と\(B\)について、
\begin{equation}
[A,B]=AB-BA.
\end{equation}
2つの観測可能量が交換するなら、それらは同時に指定できる。交換しないなら、量子力学はそれらの同時の鋭い指定を妨げる。

3つの位置座標はすべて一緒に測定できるので、位置成分は交換する：
\begin{equation}
[x_i,x_j]=0.
\end{equation}
特に、すべての演算子はそれ自身と交換する。ある観測可能量が、その同じ観測可能量の別の測定と両立しないなどということは意味をなさない。

運動量成分も同様に交換する：
\begin{equation}
[p_i,p_j]=0.
\end{equation}
位置成分は異なる方向の運動量成分とは交換するが、対応する運動量成分とは交換しない。3つの正準な命題はまとめて
\begin{equation}
[x_i,x_j]=0,\qquad
[p_i,p_j]=0,\qquad
[x_i,p_j]=i\hbar\,\delta_{ij}.
\end{equation}
と書ける。したがって、例えば
\[
[x,p_y]=0,\qquad [x,p_x]=i\hbar,
\]
であり、順序を逆にすると符号が逆になる：
\begin{equation}
[p_j,x_i]=-i\hbar\,\delta_{ij}.
\end{equation}

角運動量とプランク定数は同じ次元をもつ。角運動量は距離×運動量であり、一方、不確定性関係は次元の形で
\begin{equation}
\Delta x\,\Delta p\sim\hbar
\end{equation}
（正確な不等式におけるおなじみの数値因子を除いて）となる。したがって、角運動量が\(\hbar\)の単位で量子化されるのは自然であるが、その量子化はまだ導かれていない。

\section{角運動量代数の導出}

最初の課題は、角運動量の2つの成分の交換子を計算することである。
\[
L_x=yp_z-zp_y,\qquad
L_y=zp_x-xp_z,
\]
から
\begin{align}
[L_x,L_y]
&=[yp_z-zp_y,\;zp_x-xp_z] \\
&=[yp_z,zp_x]-[yp_z,xp_z] \notag\\
&\quad{}-[zp_y,zp_x]+[zp_y,xp_z].
\end{align}
が得られる。

\begin{classroomqa}
\lecturetimestamp{00:31:41}
\audiencequestion{3つの式を上から下へ巡回させる代わりに、記号を左から右へ巡回させることで成分の公式を覚えられるか。}
\lecturerresponse{はい。どちらの巡回の覚え方でもよい。この覚え方の唯一の目的は、符号を追いながら正しい巡回順序で成分を生成することである。}
\end{classroomqa}

交換子は、一方の積の因子をもう一方の積の因子の間を通過させられるかどうかを問うことで理解できる。
\[
[yp_z,xp_z],
\]
において、演算子\(y\)は\(x\)とも\(p_z\)とも可換であり、\(p_z\)は\(x\)ともそれ自身とも可換である。障害はないので、この項は消える。同じ推論により
\[
[zp_y,zp_x]=0.
\]
が得られる。生き残る積の組は2つだけである。

最初のものについては、唯一の障害は\(p_z\)が\(z\)を通過することである：
\begin{align}
[yp_z,zp_x]
&=yp_zzp_x-zp_xyp_z \\
&=y(p_zz-zp_z)p_x \\
&=y[p_z,z]p_x \\
&=-i\hbar\,yp_x.
\end{align}
正準交換子の順序が重要である：
\[
[p_z,z]=-[z,p_z]=-i\hbar.
\]

もう一方の生き残る組については、\(L_x\)と\(L_y\)の定義における負号が掛け合わさって正号になる。唯一の障害は\(z\)が\(p_z\)を通過することである：
\begin{align}
[zp_y,xp_z]
&=zp_yxp_z-xp_zzp_y \\
&=x(zp_z-p_zz)p_y \\
&=x[z,p_z]p_y \\
&=i\hbar\,xp_y.
\end{align}
2つの寄与を組み合わせると、
\begin{align}
[L_x,L_y]
&=-i\hbar\,yp_x+i\hbar\,xp_y \\
&=i\hbar(xp_y-yp_x) \\
&=i\hbar L_z.
\end{align}

これが重要な単純化である。角運動量の2つの成分を交換すると第3の成分が生じ、位置や運動量の演算子は残らない。巡回置換により
\begin{equation}
\begin{aligned}
[L_x,L_y]&=i\hbar L_z,\qquad
[L_y,L_z]=i\hbar L_x,\\
[L_z,L_x]&=i\hbar L_y.
\end{aligned}
\end{equation}
が得られる。等価な形では、
\begin{equation}
[L_i,L_j]=i\hbar\,\epsilon_{ijk}L_k.
\end{equation}

代数は角運動量演算子自身の上で閉じている。これらの関係が得られれば、それらを導くのに用いた\(x_i\)と\(p_i\)は脇に置いてよい。同じ代数は、全角運動量が多数の構成要素にわたる和である場合にも当てはまる。特に、残る角運動量がスピンである重心系で当てはまる。したがって、角運動量の交換子は、スピンを小さな構成部分の運動としてまず構成することなく、抽象的に用いることができる。

\(x\)、\(y\)、\(z\)軸の選び方は任意である。他の任意の互いに垂直な慣性デカルト軸に回転させても、プライム付きの成分は同じ関係に従う：
\begin{equation}
[L'_x,L'_y]=i\hbar L'_z,
\end{equation}
およびその巡回的な相方も同様である。代数は回転不変であるから、物理的予測は計算にどの慣性軸を選んだかに依存しえない。

\begin{classroomqa}
\lecturetimestamp{00:39:09}
\audiencequestion{スピン角運動量は、座標軸の回転だけでなく並進にも独立か。}
\lecturerresponse{はい。固有のスピンは並進に独立である。原点のまわりの軌道角運動量は原点を並進させると変化しうるが、重心系に残るスピンはその並進に依存しない。}
\end{classroomqa}

抽象代数は無限に多くの恒等式を生成しうるが、そのほとんどは物理的に興味深くない。その価値は、恒等式が測定可能な事実になったり、測定の操作的意味について何かを明らかにしたりするときに現れる。ここまでは構成はまだ主に代数的である。次の段階ではスピンの離散構造を引き出す。

\begin{classroomqa}
\lecturetimestamp{00:40:42}
\audiencequestion{選んだ\(x\)、\(y\)、\(z\)軸が粒子とともに回転したらどうなるか。}
\lecturerresponse{議論は慣性系に限られる。回転座標系は非慣性系であり、遠心力やコリオリ力のような効果を導入する。角運動量は共回転系への変更の下で不変ではない：慣性系で明らかに回転している系が、それとともに回転する系では静止して見えうる。ここで用いられる回転不変性は、向きの異なる慣性デカルト軸に関するものであり、時間とともに回転する軸に関するものではない。}
\end{classroomqa}

\section{昇降演算子と量子化軸}

2つの組み合わせ
\begin{equation}
L_+=L_x+iL_y,
\qquad
L_-=L_x-iL_y.
\end{equation}
を導入する。これらは調和振動子の生成演算子と消滅演算子に類似している。その役割は、選んだ角運動量の1つの成分を上げ下げすることである。

\(z\)軸に特別なものはないが、成分を指定するには1つの軸を選ばなければならない。\(z\)を量子化軸として
\begin{equation}
L_z\lvert m\rangle=m\hbar\,\lvert m\rangle.
\end{equation}
と書く。したがって\(m\)は、\(\hbar\)を単位として測定された\(z\)成分である。代数の途中では\(\hbar=1\)と置くと便利なことが多いが、ここでは明示的に保持する。

交換子\([L_+,L_-]\)は\(L_z\)に比例するが、当面の議論には必要ない。まず重要なのは、各昇降演算子と\(L_z\)との交換子である。
\[
[L_x,L_z]=-i\hbar L_y,
\qquad
[L_y,L_z]=i\hbar L_x,
\]
から
\begin{align}
[L_+,L_z]
&=[L_x+iL_y,L_z] \\
&=[L_x,L_z]+i[L_y,L_z] \\
&=-i\hbar L_y+i(i\hbar L_x) \\
&=-\hbar(L_x+iL_y) \\
&=-\hbar L_+.
\end{align}
が得られる。\(iL_y\)の項の符号を変えると
\begin{equation}
[L_-,L_z]=+\hbar L_-.
\end{equation}
となる。等価な表現として、各交換子の順序を逆にすれば
\begin{equation}
[L_z,L_+]=\hbar L_+,
\qquad
[L_z,L_-]=-\hbar L_-.
\end{equation}

角運動量の成分のうち、同時に確定できるのは一つだけである。異なる二つの成分は交換しないため、同時に指定することはできない。選ばれた確定成分が量子化軸を定義する。\(z\)を選ぶのは便宜上のことであり、物理的な選好ではない。

次に、\(L_+\)を\(L_z\)の固有状態に作用させる。交換子は
\begin{equation}
L_zL_+=L_+L_z+\hbar L_+.
\end{equation}
と書き直せる。したがって
\begin{align}
L_z\bigl(L_+\lvert m\rangle\bigr)
&=\bigl(L_+L_z+\hbar L_+\bigr)\lvert m\rangle \\
&=m\hbar L_+\lvert m\rangle+\hbar L_+\lvert m\rangle \\
&=(m+1)\hbar\,L_+\lvert m\rangle.
\end{align}
となる。ゼロでない限り、\(L_+\lvert m\rangle\)は\(L_z\)の固有状態であり、その固有値は\(\hbar\)の1単位分だけ増加している。同様に
\begin{equation}
L_z\bigl(L_-\lvert m\rangle\bigr)
=(m-1)\hbar\,L_-\lvert m\rangle.
\end{equation}
であるから、\(L_-\)は\(m\)を1だけ下げる。

これはすでに自明でない結論である。\(L_z\)のスペクトルがどうであれ、隣り合う値は整数で隔てられている：
\begin{equation}
\ldots,\ m-2,\ m-1,\ m,\ m+1,\ m+2,\ \ldots
\end{equation}
端点はまだ決定されていない。

\section{有限個の多重項とスピンラベル}

一定の種類の粒子は一定の角運動量の大きさをもつ。その角運動量は向きを変えうるが、大きさを際限なく増やせば、やがて別の対象が生じることになる。したがって、その粒子には最大可能な\(z\)成分が存在しなければならない。

昇降演算子をもつ梯子に上限がありうるのはなぜか。調和振動子がその模型を与える。その消滅演算子は最下状態に達するとゼロを与える。ここでは、昇降演算子は最高状態に達するとゼロを与えなければならない：
\begin{equation}
L_+\lvert m_{\max}\rangle=0.
\end{equation}
\(L_+\)が固有値を上げるという主張は、それが生み出すベクトルがゼロでない場合にのみ当てはまる。最高ウェイト状態とは、まさに昇降列が終止する例外状態である。

回転不変性が下端を固定する。角運動量が\(+z\)に沿って最大限に向くことができるなら、同じ物理的対象を回転させて\(-z\)に沿って最大限に向くようにできる。したがって
\begin{equation}
m_{\min}=-m_{\max}.
\end{equation}
スペクトルはゼロに関して対称であり、間隔は一単位である。可能な系列は二つある。

第一の系列では、ゼロが現れる：
\begin{equation}
\begin{aligned}
m={}&-m_{\max},-m_{\max}+1,\ldots,-1,0,\\
&1,\ldots,m_{\max}-1,m_{\max}.
\end{aligned}
\end{equation}
このとき\(m_{\max}\)は整数であり、状態の数は
\begin{equation}
2m_{\max}+1.
\end{equation}
である。最大値がその粒子を特徴づけ、そのスピンと呼ばれる。慣用的な表記では、
\begin{equation}
s\equiv m_{\max},
\qquad
\lvert s,m\rangle\equiv\lvert m\rangle
\end{equation}
であり、固定スピンの多重項内で成り立つ。\footnote{\raggedright 講義動画フレーム\nolinkurl{lecture_07_figure_03.png}、01：01：08。\par}

\begin{figure}[htbp]
\centering
\includegraphics[width=0.70\textwidth]{lecture_07_figure_03.png}
\caption{最高ウェイト条件が昇降列を終止させ、許される\(m\)の値は有限個の単位間隔の梯子をなす。}
\lectureframe{01:01:08}
\end{figure}

スピンゼロの粒子は\(m=0\)のみをもつ。スピン1の粒子は
\[
m=1,0,-1,
\]
をもち、スピン2の粒子は五つの成分値をもつ。ヒッグス粒子はスピンゼロの例である。光子と\(Z\)ボソンはスピン1をもつ。重力子が存在するならばスピン2をもつ。複合的な対象はスピン3やそれ以上のはるかに大きなスピンをもちうる。原子核はそうでありうるし、バスケットボールや銀河のような十分に巨視的な回転物体は任意に大きな角運動量を担いうる。通常の素粒子の一覧に欠けているのは、素粒子としてのスピン3の粒子である。

もう一方の対称な単位間隔の系列はゼロを含まない。それは
\[
m=\frac12,-\frac12
\]
で始まり、単位刻みで広がる：
\begin{equation}
\begin{aligned}
m={}&-m_{\max},-m_{\max}+1,\ldots,-\frac32,-\frac12,\\
&\frac12,\frac32,\ldots,m_{\max}.
\end{aligned}
\end{equation}
ここで\(m_{\max}\)は半整数である。最も単純な成員はスピン\(\tfrac12\)をもち、続いてスピン\(\tfrac32\)、スピン\(\tfrac52\)、…となる。したがって角運動量の多重項は整数スピン系列と半整数スピン系列に分かれる。

\begin{classroomqa}
\lecturetimestamp{01:03:46}
\audiencequestion{半整数スピンの多重項にも最大値はあるのか。}
\lecturerresponse{はい。梯子にはやはり最高値がある。スピン\(\tfrac32\)の粒子では最大値は\(m_{\max}=\tfrac32\)であり、スピン\(\tfrac52\)の粒子では\(m_{\max}=\tfrac52\)である。違いは、最大値および多重項内の他のすべての値が半整数であるという点だけである。}
\end{classroomqa}

\begin{classroomqa}
\lecturetimestamp{01:04:04}
\audiencequestion{スピン3以上の素粒子は存在しないのか。}
\lecturerresponse{ここで議論している相対論的量子論の枠組みでは、通常の一覧にある素粒子のスピンは\(0,\tfrac12,1,\tfrac32,\)および\(2\)である。複合的な対象はこの範囲に制限されず、スピン3あるいはそれよりはるかに大きなスピンをもちうる。}
\end{classroomqa}

\begin{classroomqa}
\lecturetimestamp{01:04:50}
\audiencequestion{はしごの負の成分は負のスピンを表すのですか。}
\lecturerresponse{いいえ。負の\(m\)は成分の値であり、選んだ軸に対して角運動量が逆向きを向いていることを示す。粒子のスピンは非負の最大値\(s=m_{\max}\)である。}
\end{classroomqa}

\begin{classroomqa}
\lecturetimestamp{01:05:10}
\audiencequestion{成分の値が離散的なら、その間のすべての向きはどうなるのですか。}
\lecturerresponse{可能な空間的な向きは二つや三つだけではない。中間の向きは\(m\)基底状態の量子重ね合わせで表される。以下のスピン\(\tfrac12\)の構成がこれを明示する。}
\end{classroomqa}

\begin{classroomqa}
\lecturetimestamp{01:05:42}
\audiencequestion{コルク抜き運動や円偏光波はどうですか。}
\lecturerresponse{円偏光した電磁波は円偏光した光子から構成され、それらの光子は伝播軸に沿った角運動量を持つ。``コルク抜き''という語は回転と前進運動を組み合わせているので、静止していると考えられる物体の記述には適切でない。運動する粒子については、関連する量はヘリシティ、すなわち運動方向に対するスピンの向きである。}
\end{classroomqa}

\begin{classroomqa}
\lecturetimestamp{01:07:10}
\audiencequestion{スピン\(\tfrac32\)の既知の素粒子はありますか。}
\lecturerresponse{スピン\(\tfrac32\)の素粒子は直接確認されていない。超対称性はグラビトンと呼ばれるスピン\(\tfrac32\)の重力子のパートナーを予想しているが、それは直接検出されていない。}
\end{classroomqa}

\begin{classroomqa}
\lecturetimestamp{01:08:01}
\audiencequestion{角運動量演算子は元の波動関数やケット空間に作用するのをやめたのですか。}
\lecturerresponse{いいえ、元のスピン空間はまだ指定されていなかった。この代数はその空間を決定するために用いられている。固定されたスピン\(s\)に対して、許される値でラベル付けされた基底ベクトルを持つ
\[
\lvert s,m\rangle,\qquad m=-s,-s+1,\ldots,s.
\]
この内在的スピン空間は、軌道角運動量が\(\mathbf r\times\mathbf p\)として作用する空間的波動関数の通常の空間とは異なるが、両者は同じ角運動量代数を実現している。}
\end{classroomqa}

\section{角運動量の大きさ}

\begin{classroomqa}
\lecturetimestamp{01:09:32}
\audiencequestion{角運動量ベクトルの大きさは観測可能ですか。}
\lecturerresponse{はい。その二乗を用いるのが最も便利である。
\[
L^2=L_x^2+L_y^2+L_z^2.
\]
最高ウェイト状態におけるその値は単に\(s^2\hbar^2\)ではない。\(L_z\)が確定して最大であっても、非可換な横成分は両方とも鋭くゼロにはなり得ないので、それらは追加の量子寄与を与える。}
\end{classroomqa}

その寄与を計算するには、上昇演算子がケットに先に作用する順序ではしご演算子を掛ける。
\begin{align}
L_-L_+
&=(L_x-iL_y)(L_x+iL_y) \\
&=L_x^2+L_y^2+iL_xL_y-iL_yL_x \\
&=L_x^2+L_y^2+i[L_x,L_y] \\
&=L_x^2+L_y^2-\hbar L_z.
\end{align}
最後のステップでは\([L_x,L_y]=i\hbar L_z\)を用いる。したがって
\begin{equation}
L_x^2+L_y^2=L_-L_++\hbar L_z,
\end{equation}
となり、ゆえに
\begin{equation}
L^2=L_-L_++L_z^2+\hbar L_z.
\end{equation}

この恒等式を最高ウェイト状態\(\lvert s,s\rangle=\lvert m_{\max}\rangle\)に適用する。
\begin{align}
L^2\lvert s,s\rangle
&=\left(L_-L_++L_z^2+\hbar L_z\right)\lvert s,s\rangle \\
&=L_-L_+\lvert s,s\rangle
+s^2\hbar^2\lvert s,s\rangle
+s\hbar^2\lvert s,s\rangle.
\end{align}
\(L_+\lvert s,s\rangle=0\)なので第一項は消える。したがって
\begin{equation}
L^2\lvert s,s\rangle
=s(s+1)\hbar^2\lvert s,s\rangle.
\end{equation}
余分な\(+s\)は避けられない横方向の揺らぎに伴う量子補正である。古典的には\(s^2\hbar^2\)のみが期待されるであろう。

最初のいくつかの値がこの区別を具体的にする。
\begin{align}
s=0 &: \qquad L^2=0,\\
s=1 &: \qquad L^2=2\hbar^2,\\
s=\frac12 &: \qquad L^2=\frac34\hbar^2.
\end{align}
\(s=1000\)に対して、無次元因子は\(1000\times1001\)であり、このような大きな角運動量の尺度では\(1000^2\)とほとんど変わらない。一般に、
\begin{equation}
s(s+1)=s^2\left(1+\frac1s\right),
\end{equation}
であるから、非常に大きな\(s\)に対して相対的な量子補正は無視できるようになる。これが古典的な\(s^2\)の結果が現れる仕方である。

固定されたスピン多重項内のすべての\(m\)状態は、\(L^2\)の同じ固有値を持つ。\(m\)を変えることは、選んだ軸に沿った成分を変えるのであり、全大きさを変えるのではない。これは、角運動量ベクトルの長さを変えずに回転させることの量子版である。

\begin{classroomqa}
\lecturetimestamp{01:16:43}
\audiencequestion{\(L^2\)と\(z\)成分を同時に知ることはできますか。}
\lecturerresponse{はい。全二乗はすべての角運動量成分と可換である。
\[
[L^2,L_i]=0,\qquad i=x,y,z.
\]
したがって、二つの異なる成分は同時に指定できないが、\(L^2\)と一つの選ばれた成分、たとえば\(L_z\)は同時に指定できる。たとえばすべての電子は、その向きによらず\(L^2=\tfrac34\hbar^2\)を持つ。}
\end{classroomqa}

\begin{classroomqa}
\lecturetimestamp{01:17:41}
\audiencequestion{\(L^2\)がゼロなのはスピン1に対してですか、それともスピンゼロに対してのみですか。}
\lecturerresponse{\(L^2=0\)となるのはスピンゼロのみである。スピン1は
\[
L^2=1(1+1)\hbar^2=2\hbar^2,
\]
を与え、一方スピン3は
\[
L^2=3(3+1)\hbar^2=12\hbar^2.
\]}
を与える。\end{classroomqa}

\begin{classroomqa}
\lecturetimestamp{01:18:13}
\audiencequestion{なぜ多重項の各\(m\)準位は同じ\(L^2\)固有値をもつのか。}
\lecturerresponse{最高重み状態から出発してそれを下げる。\(L^2\)はすべての成分と可換であるから、\(L_-\)とも可換である。
\[
[L^2,L_-]=0.
\]
もし
\[
L^2\lvert s,s\rangle=s(s+1)\hbar^2\lvert s,s\rangle,
\]
ならば
\begin{align}
L^2\bigl(L_-\lvert s,s\rangle\bigr)
&=L_-L^2\lvert s,s\rangle \\
&=s(s+1)\hbar^2\,L_-\lvert s,s\rangle.
\end{align}
したがって、次に下がった状態も同じ\(L^2\)固有値をもつ。この議論を繰り返すと、その固有値ははしごの一番下まで引き継がれる。}
\end{classroomqa}

したがって、\(m\)の値が離散的であることは、空間的な向きが離散的な集合しか存在しないことを意味しない。回転不変性は依然として、回転させた向きに対応する状態を要求する。スピン\(\tfrac12\)の場合、それらの向きは2つの\(L_z\)基底状態の重ね合わせとして現れる。

\section{スピン\texorpdfstring{\(\tfrac12\)}{1/2}の状態と向き}

回転不変性は離散的な\(L_z\)固有状態だけでは足りない。スピンは任意の軸に沿って向きうるが、選んだ軸に沿った成分の測定は許された離散値しか与えない。\(z\)軸固有状態の間の向きは重ね合わせで表される。

スピン\(\tfrac12\)の粒子について、スピン自由度だけに注目する。電子には軌道状態もあるので、全体として2つの状態しかもつわけでは決してない。そのスピン空間だけが2次元である。\(z\)軸基底を選ぶと
\begin{equation}
L_z\lvert+\rangle=\frac{\hbar}{2}\lvert+\rangle,
\qquad
L_z\lvert-\rangle=-\frac{\hbar}{2}\lvert-\rangle.
\end{equation}
一般のスピン状態は
\begin{equation}
\lvert\psi\rangle
=
\alpha\lvert+\rangle+\beta\lvert-\rangle,
\end{equation}
ここで\(\alpha\)と\(\beta\)は複素数である。\(L_z\)の測定は確率\(|\alpha|^2\)で\(+\hbar/2\)を、確率\(|\beta|^2\)で\(-\hbar/2\)を与えるので、規格化には
\begin{equation}
|\alpha|^2+|\beta|^2=1.
\end{equation}
が必要である。

共通の位相は物理的効果を持たない：
\begin{equation}
(\alpha,\beta)\longrightarrow
e^{i\chi}(\alpha,\beta).
\end{equation}
二つの複素振幅は四つの実パラメータから始まる。規格化により一つが除かれ、無関係な全体位相によりさらに一つが除かれ、二つの物理的実パラメータが残る。これはまさに、極角と方位角によって方向を指定するのに必要な数である。\(\alpha\) と \(\beta\) を変化させると、スピンの向きが回転する。

測定された成分とその期待値は別物である。期待値はアンサンブル平均である。多数の電子をまったく同じ状態に準備し、各電子について\(L_z\)を1回測定し、結果を足して実験回数で割る。個々の結果は依然として\(+\hbar/2\)か\(-\hbar/2\)のどちらかであるが、平均はそれらの間の任意の値をとりうる。
\begin{equation}
\langle L_z\rangle
=
\frac{\hbar}{2}\bigl(|\alpha|^2-|\beta|^2\bigr).
\end{equation}

上向きと下向きの振幅が等しいと、スピンは\(x\)軸の向きの一方に沿って向く。規格化された代表は\footnote{係数の比が向きを定める。因子\(1/\sqrt2\)は\( |\alpha|^2+|\beta|^2=1\)が要求する規格化を与える。相対位相\(+i\)と\(-i\)は逆の\(y\)軸の向きを与える。正の\(y\)にどちらの符号を割り当てるかは軸の規約に依存する。}
\begin{equation}
\lvert x;+\rangle
=
\frac{\lvert+\rangle+\lvert-\rangle}{\sqrt2},
\qquad
\lvert x;-\rangle
=
\frac{\lvert+\rangle-\lvert-\rangle}{\sqrt2}.
\end{equation}
どちらの状態でも、\(z\)に沿った測定は2つの可能な結果を等しい確率で与えるので、
\begin{equation}
\langle L_z\rangle
=
\frac12\left(\frac{\hbar}{2}\right)
+
\frac12\left(-\frac{\hbar}{2}\right)
=0.
\end{equation}
プラスとマイナスの相対符号は\(x\)軸に沿った逆向きを与える。同様に、
\begin{equation}
\frac{\lvert+\rangle+i\lvert-\rangle}{\sqrt2},
\qquad
\frac{\lvert+\rangle-i\lvert-\rangle}{\sqrt2}
\end{equation}
は\(y\)軸の逆向きの2つの向きを与える。\(\alpha\)と\(\beta\)の一般の規格化された値は任意の軸を表す。スピンがある特定の軸に沿って向くということは、その軸に沿った成分の測定が確定した結果を与えることを意味する。別の軸に対しては、同じ状態は一般に重ね合わせである。

個々の測定結果は量子化されたままである。連続的に変化し、古典的なベクトルのおなじみの幾何学を再現するのは、それらのアンサンブル平均である。

\section{スピン、統計、複合粒子}

質量、電荷、スピンは粒子を特徴づける主要な性質の一部である。スピンは統計とも相関する。スピン統計関係は相対論的量子力学の定理であるが、ここでは証明しない。
\begin{equation}
\begin{aligned}
\text{half-integer spin} &\quad\Longleftrightarrow\quad \text{fermion},\\
\text{integer spin} &\quad\Longleftrightarrow\quad \text{boson}.
\end{aligned}
\end{equation}
2つの同種フェルミオンは同じ完全な量子状態を占めることができない。同種ボソンは同じ状態を占めることができ、多数が1つの状態にあると古典的な波のように集団的に振る舞いうる。したがって、スピン\(\tfrac12,\tfrac32,\tfrac52,\ldots\)の粒子はフェルミオンであり、スピン\(0,1,2,\ldots\)の粒子はボソンである。

この規則は直ちに複合粒子に拡張される。2つの半整数角運動量を加えると整数角運動量になる。より一般に、偶数個のフェルミオン的構成要素はボソン的複合粒子をつくり、奇数個はフェルミオン的複合粒子をつくる。ボソン的構成要素を加えてもこの偶奇規則は変わらない。

通常の水素はスピン\(\tfrac12\)の陽子とスピン\(\tfrac12\)の電子を含む。したがって偶数個のフェルミオンを含み、全体としてボソン的である。重陽子は陽子と中性子を含み、どちらもフェルミオンであり、重陽子もまたボソンである。重水素原子は重陽子に電子を加えたもので、この記述のレベルでは三つのフェルミオン的構成要素を与えるので、重水素はフェルミオンである。

粒子と反粒子は同じ質量とスピンをもち、電荷は反対である。ポジトロニウムは電子と陽電子が互いに束縛されたものである。二つのフェルミオンを含むので、ボソン的である。

\begin{classroomqa}
\lecturetimestamp{01:33:08}
\audiencequestion{ボソンは物を通り抜けるのか、また超伝導は対になった電子がボソンとして振る舞うことと関係しているのか。}
\lecturerresponse{関係する性質は、単に物質を通り抜ける能力ではなく、ボソンが同じ状態を占めることができるという性質である。超流動ヘリウムがより簡単な例を与える。通常のヘリウム\(4\)原子は二つの陽子、二つの中性子、二つの電子を含み、偶数個のフェルミオン的構成要素をもつので、この原子はボソンである。多くのヘリウム原子が一つの量子状態を占め、集団的に運動することができる。すると超流体は、多くの原子を同じ状態に積み重ねてできる古典波として、近似的に描くことができる。その描像は厳密ではないが、ここで必要な要点を捉えている。金属中の電子はフェルミオンであり、別の電子と完全な状態を共有することはできない。しかし電子の対はボソン的複合体であり、これが超伝導とのつながりである。ここで逆説が生じる。構成要素である電子が同じ状態を占めることができないなら、いくつかの複合対が一つの状態を占めることはどうしてできるのか。その解決には複合状態についてより注意深い説明が必要であり、後回しにする。}
\end{classroomqa}

\begin{classroomqa}
\lecturetimestamp{01:35:46}
\audiencequestion{陽子と中性子は複合粒子であるのに、なぜ半整数スピンをもつのか。}
\lecturerresponse{それぞれは奇数個のスピン\(\tfrac12\)クォークからできている。三つのフェルミオン的構成要素がフェルミオン的複合体をつくるので、陽子と中性子は半整数スピンをもつことができる。}
\end{classroomqa}

\section{パウリ排他律と原子軌道}

周期表はパウリ排他原理の直接の動機を与える。規則は、二つの電子が決して同じ軌道運動をもてないということではない。二つの電子が同じ完全な量子状態を占めることはできないということである。

この区別は最も単純な原子でも重要である。水素から始めて、核電荷を増やし、第二の電子を加える。二つの電子は同じ空間軌道を占めてもよいが、それらのスピン状態は異ならなければならない。選んだ軸に対して一方がスピン上向き、他方がスピン下向きであってよい。二つの電子が一つの軌道を占めるとき、それらのスピンはスピン一重項状態に結合するので、全スピンは打ち消し合う。利用できる独立なスピン状態は二つしかないので、一つの空間軌道を占めることができる電子は二つまでである。

\begin{classroomqa}
\lecturetimestamp{01:37:53}
\audiencequestion{一つの軌道にある二つの電子の反対向きのスピンは、かみ合った歯車として描けるか。}
\lecturerresponse{粗い描像としては、はい。かみ合った歯車は反対向きに回る。それは反対向きのスピン方向のための類推にすぎず、電子の文字通りの力学模型ではない。}
\end{classroomqa}

``軌道''という語は二つの異なるレベルで用いることができる。広い意味での軌道運動の状態とは、スピンを無視してシュレーディンガー方程式を解くことによって得られる状態である。それは電子がどの殻を占めるかだけよりも多くの情報を含む。

\begin{classroomqa}
\lecturetimestamp{01:38:09}
\audiencequestion{軌道とは、電子がどの殻を占めるかだけ以上のものを意味するのか。}
\lecturerresponse{はい。ここで軌道運動とは、電子のスピンを無視したシュレーディンガー状態を意味する。それは動径情報を担い、電子の原子核からの距離をおよそ定める主量子数\(n\)、および軌道角運動量量子数を含む。化学者が用いる分光表記では、\(S\)、\(P\)、\(D\)、\(F\)はそれぞれ軌道角運動量\(l=0,1,2,3\)に対応する。}
\end{classroomqa}

したがって主量子数と軌道角運動量は運動の異なる側面を記述する。\(S,P,D,F\)という記号は\(l\)の値を指し、一方\(n\)は別個の動径情報を与える。パウリの規則は、スピン量子数を含め、二つの電子がそれらの量子数のすべてを同じにすることはできないというものである。

\begin{classroomqa}
\lecturetimestamp{01:40:42}
\audiencequestion{量子数によって、原子の各殻に何個の電子が入り得るかを決めることができるか。}
\lecturerresponse{はい。量子数は可能な最大数を決めるが、特定の原子がその殻に実際に何個の電子を含むかは決めない。一般的な殻の数え上げの問題はここでは扱わない。}
\end{classroomqa}

固定された軌道角運動量\(l\)に対して、成分の量子数は
\begin{equation}
2l+1
\end{equation}
個の値をとり得る。これはスピンですでに用いた角運動量の数え上げと同じである。\(S\)状態は\(l=0\)なので、1個の軌道状態をもつ。\(P\)状態は\(l=1\)なので、3個の軌道状態をもつ。\(D\)状態は\(l=2\)なので、5個をもつ。

スピンは各軌道の数を2倍にする。なぜなら、各空間軌道は、互いに逆向きの2つのスピン状態のそれぞれに1個ずつ電子を収容できるからである。

\begin{classroomqa}
\lecturetimestamp{01:41:54}
\audiencequestion{\(S\)状態は2個の電子を収容し、3つの\(P\)状態はさらに6個を収容するのか。}
\lecturerresponse{はい。単一の\(S\)軌道状態は最大2個の電子を収容する。3つの\(P\)軌道状態はそれぞれ2個ずつを収容し、さらに6個を与える。したがって、ここで議論している\(S\)状態と\(P\)状態は、合計\(2+6=8\)個の電子を収容する。}
\end{classroomqa}

原子軌道を組織するのと同じ角運動量の構造が、いくつかの違いを伴いつつ、クォークが束縛系の中で運動し互いの周りを回るときにも再び現れる。
